% Transcription of Cayley's Collected Mathematical Papers, Vol. XII
% PDF pages 226--237 of collmathpapers12caylrich.pdf
% (book pages 206--217: Paper 825 continued -- Memoir on the Abelian and Theta
%  Functions, Ch. VI conclusion (Arts 162--163), Ch. VII (Arts 164--184);
%  and Paper 826 -- Note on a Partition-Series, opening pages.)
% Source scan: local workspace\Documents\Papors\OS\Cayley\collmathpapers12caylrich.pdf
\documentclass[11pt]{article}
\usepackage[a4paper,margin=0.65in]{geometry}
\usepackage[T1]{fontenc}
\usepackage{amsmath,amssymb,amsfonts}
\usepackage{graphicx}
\usepackage{pdfpages}
\usepackage[english,shorthands=off]{babel}
\usepackage{fancyhdr}
\usepackage[bookmarks=false,colorlinks=false]{hyperref}
\usepackage[protrusion=true,expansion=false]{microtype}
\emergencystretch=3em
\tolerance=600
\providecommand{\modd}{\mathrm{mod}}
\providecommand{\Disc}{\mathrm{Disc}}
\providecommand{\canont}{}
\providecommand{\asterism}{*}
\providecommand{\Hessian}{H}
\providecommand{\tome}{}
\providecommand{\page}{p.}
\pagestyle{fancy}\fancyhf{}\fancyhead[L]{Pages 226--237}\fancyhead[R]{Cayley --- Collected Papers, Vol. XII}\renewcommand{\headrulewidth}{0.4pt}
\setlength{\headheight}{15pt}

\begin{document}

\begin{center}
{\large\scshape The Collected Mathematical Papers of Arthur Cayley}\\[0.3em]
{\large Vol.\ XII\quad (Scan PDF pp.\ 226--237; book pp.\ 206--217)}\\[0.4em]
{\itshape Paper 825 (Memoir on the Abelian and Theta Functions, concl.\ of Ch.\ VI;
Ch.\ VII---The Functions $T$, $U$, $V$, $\Theta$); Paper 826 (Note on a Partition-Series).}
\end{center}
\vspace{0.6em}

% =====================================================================
% Page 206 of book (scan p. 226)  -- Paper 825 cont., Art. 161 end + 162
% =====================================================================

\noindent\textsc{206 \hfill A Memoir on the Abelian and Theta Functions. \hfill [825}
\vspace{0.6em}

\noindent
If, reverting to the identity, we write therein for instance $x=a$, we find
\[
   \alpha a^{3}+\beta a^{2}+\gamma a+\delta = \sqrt{\mu}\,A_{a}A_{m}A_{m},
\]
which equation when properly reduced gives the proportional value of $A_{m}$.

\medskip
\textbf{162.}\quad Calling for a moment the function on the left-hand side $\Omega$,
we have
\[
\Omega
\;=\;
\begin{vmatrix}
 x_{1}^{3} & x_{1}^{2} & x_{1} & 1 & \sqrt{\lambda}\sqrt{X_{1}}\\[2pt]
 x_{2}^{3} & x_{2}^{2} & x_{2} & 1 & \sqrt{\lambda}\sqrt{X_{2}}\\[2pt]
 x_{3}^{3} & x_{3}^{2} & x_{3} & 1 & \sqrt{\lambda}\sqrt{X_{3}}\\[2pt]
 x_{4}^{3} & x_{4}^{2} & x_{4} & 1 & \sqrt{\lambda}\sqrt{X_{4}}\\[2pt]
 a^{3} & a^{2} & a & 1 & \Omega
\end{vmatrix}
\;=\;0,
\]
that is,
\[
\resizebox{\textwidth}{!}{$\displaystyle
\Omega
\begin{vmatrix} x_{1}^{2} & x_{1} & 1\\ x_{2}^{2} & x_{2} & 1\\ x_{3}^{2} & x_{3} & 1\\ x_{4}^{2} & x_{4} & 1\end{vmatrix}
\;+\;\sqrt{\lambda}
\begin{vmatrix} x_{1}^{3} & x_{1}^{2} & x_{1} & \sqrt{X_{1}}\\ x_{2}^{3} & x_{2}^{2} & x_{2} & \sqrt{X_{2}}\\
                x_{3}^{3} & x_{3}^{2} & x_{3} & \sqrt{X_{3}}\\ x_{4}^{3} & x_{4}^{2} & x_{4} & \sqrt{X_{4}}\end{vmatrix}=0,
$}
\]
viz.\ this is
\[
\Omega\,(x_{1}-x_{2})(x_{1}-x_{3})(x_{1}-x_{4})(x_{2}-x_{3})(x_{2}-x_{4})(x_{3}-x_{4})
\]
\[
\resizebox{\textwidth}{!}{$\displaystyle
\;=\;-\sqrt{\lambda}\,\bigl\{
 \sqrt{X_{1}}\cdot x_{2}-x_{3}\cdot x_{2}-x_{4}\cdot x_{3}-x_{4}\cdot a-x_{1}\cdot x_{2}-a\cdot x_{3}-a\cdot x_{4}-a
$}
\]
\[
\resizebox{\textwidth}{!}{$\displaystyle
\quad+\sqrt{X_{2}}\cdot x_{1}-x_{3}\cdot x_{1}-x_{4}\cdot x_{3}-x_{4}\cdot a-x_{2}\cdot x_{1}-a\cdot x_{3}-a\cdot x_{4}-a
$}
\]
\[
\resizebox{\textwidth}{!}{$\displaystyle
\quad+\sqrt{X_{3}}\cdot x_{1}-x_{2}\cdot x_{1}-x_{4}\cdot x_{2}-x_{4}\cdot a-x_{3}\cdot x_{1}-a\cdot x_{2}-a\cdot x_{4}-a
$}
\]
\[
\resizebox{\textwidth}{!}{$\displaystyle
\quad+\sqrt{X_{4}}\cdot x_{1}-x_{2}\cdot x_{1}-x_{3}\cdot x_{2}-x_{3}\cdot a-x_{4}\cdot x_{1}-a\cdot x_{2}-a\cdot x_{3}-a\bigr\};
$}
\]
or, as this may be written,
\[
\Omega \;.\; x_{1}-x_{2}\;.\; x_{1}-x_{3}\;.\; x_{2}-x_{3}\;.\; x_{1}-x_{4}\;.\; x_{2}-x_{4}\;.\; x_{3}-x_{4}
\]
\[
\resizebox{\textwidth}{!}{$\displaystyle
\;=\;\frac{\sqrt{\lambda}\,.\,a-x_{1}\,.\,a-x_{2}}{x_{1}-x_{2}}\bigl[x_{3}-x_{2}\,.\,x_{3}-x_{4}\,.\,a-x_{3}\,.\,\sqrt{X_{4}}-(x_{4}-x_{2}\,.\,x_{4}-x_{3}\,.\,a-x_{4}\,.\,\sqrt{X_{3}})\bigr]
$}
\]
\[
\resizebox{\textwidth}{!}{$\displaystyle
\quad +\frac{\sqrt{\lambda}\,.\,a-x_{1}\,.\,a-x_{4}}{x_{2}-x_{4}}\bigl[x_{2}-x_{1}\,.\,x_{2}-x_{3}\,.\,a-x_{2}\,.\,\sqrt{X_{3}}-(x_{3}-x_{1}\,.\,x_{3}-x_{2}\,.\,a-x_{3}\,.\,\sqrt{X_{2}})\bigr].
$}
\]
We have here $\sqrt{\lambda}\,.\,a-x_{1}\,.\,a-x_{2}=\sqrt{\lambda}\,A_{a}^{2}A_{m}^{2}$, and the function
\[
\frac{1}{x_{3}-x_{4}}\bigl[x_{2}-x_{3}\,.\,x_{2}-x_{4}\,.\,b_{0}\sqrt{X_{4}}-(x_{1}-x_{3}\,.\,x_{1}-x_{4}\,.\,b_{0}\sqrt{X_{3}})\bigr],
\]
which multiplies this, is without difficulty found to be
\[
=\;\frac{-A_{12}}{c-d\,.\,d-b\,.\,b-c}\,\Sigma\,[c-d\,.\,B_{12}^{x}BE_{a}\,.\,C_{a}\,.\,D_{a}],
\]

% =====================================================================
% Page 207 of book (scan p. 227)
% =====================================================================
\newpage
\noindent\textsc{825] \hfill A Memoir on the Abelian and Theta Functions. \hfill 207}
\vspace{0.6em}

\noindent
where the summation extends to the three terms obtained by the cyclical interchange
of the letters $b$, $c$, $d$: these being a set of three out of the five letters other than $a$.
Similarly $\sqrt{\lambda}\,.\,a-x_{1}\,.\,a-x_{2}=\sqrt{\lambda}\,A_{a}^{2}A_{m}^{2}$, and the function which multiplies this is
\[
=\;\frac{-A_{12}}{c-d\,.\,d-b\,.\,b-c}\,\Sigma\,[c-d\,.\,B_{12}^{x}\,.\,BE_{a}C_{a}D_{a}].
\]

The expression for $\Omega$ thus contains the factor $A_{12}A_{34}$: but we have
\[
\Omega = \alpha a^{3}+\beta a^{2}+\gamma a+\delta = \sqrt{\mu}\,A_{a}A_{m}A_{m};
\]
this equation contains therefore the factor $A_{12}A_{34}$, and omitting it we find
\[
\resizebox{\textwidth}{!}{$\displaystyle
-\frac{\sqrt{\mu}}{\sqrt{\lambda}}\,(x_{1}-x_{2}\,.\,x_{1}-x_{3}\,.\,x_{1}-x_{4}\,.\,x_{2}-x_{3}\,.\,x_{2}-x_{4}\,.\,x_{3}-x_{4})\,(c-d\,.\,d-b\,.\,b-c)\,A_{m}
$}
\]
\[
\resizebox{\textwidth}{!}{$\displaystyle
\quad =\;A_{m}\,\Sigma\,[c-d\,.\,B_{12}^{x}\,BE_{a}^{1}C_{a}\,D_{a}]\;+\;A_{12}\,\Sigma\,[c-d\,.\,B_{12}^{x}\,BE_{a}^{2}C_{a}\,D_{a}],
$}
\]
where, as before, the summations refer each to the three terms obtained by the
cyclical interchange of the letters $b$, $c$, $d$; these being any three of the five letters
other than $a$: and the remaining two letters $e$, $f$ enter into the formulae symmetrically.
The formula thus gives for $A_{m}$ ten values which are of course equal to each other.

By reason of the undetermined factor $\dfrac{\sqrt{\mu}}{\sqrt{\lambda}}$, the formula gives only the proportional
value of $A_{m}$; viz.\ combining it with the like formulae for $B_{m}$, \&c., we have determinate
values of the ratios $A_{m}:B_{m}:\ldots:F_{56}$. But this being understood, we regard the
formula as a formula for each single-letter function of $x_{3}$, $x_{4}$ in terms of the single
and double-letter functions of $x_{1}$, $x_{2}$ and of $x_{3}$, $x_{4}$ respectively.

\medskip
\textbf{163.}\quad We require further the expressions for the double-letter functions of $x_{3}$, $x_{4}$.

Consider for example the function $DE_{m}$, which is
\[
=\;\frac{1}{x_{3}-x_{4}}\bigl[\sqrt{d_{3}\,e_{3}\,f_{3}\,a_{4}\,b_{4}\,c_{4}}-\sqrt{d_{4}\,e_{4}\,f_{4}\,a_{3}\,b_{3}\,c_{3}}\bigr];
\]
then multiplying by $A_{m}\,B_{m}\,C_{m}=\sqrt{a_{3}\,b_{3}\,c_{3}\,a_{4}\,b_{4}\,c_{4}}$, we have
\[
DE_{m}\,A_{m}\,B_{m}\,C_{m} \;=\; \frac{1}{x_{3}-x_{4}}\bigl[a_{3}b_{3}c_{3}\sqrt{X_{4}}-a_{4}b_{4}c_{4}\sqrt{X_{3}}\bigr],
\]
or recollecting that $\sqrt{\lambda}\,\sqrt{X_{3}}$ and $\sqrt{\lambda}\,\sqrt{X_{4}}$ are $=\alpha x_{3}^{3}+\beta x_{3}^{2}+\gamma x_{3}+\delta$ and $\alpha x_{4}^{3}+\beta x_{4}^{2}+\gamma x_{4}+\delta$
respectively, this may be written
\[
\resizebox{\textwidth}{!}{$\displaystyle
\sqrt{\lambda}\,DE_{m}\,A_{m}B_{m}C_{m} \;=\; \frac{1}{x_{3}-x_{4}}\bigl\{(a-x_{3}\,.\,b-x_{3}\,.\,c-x_{3})(\alpha x_{4}^{3}+\beta x_{4}^{2}+\gamma x_{4}+\delta)
$}
\]
\[
-(a-x_{4}\,.\,b-x_{4}\,.\,c-x_{4})(\alpha x_{3}^{3}+\beta x_{3}^{2}+\gamma x_{3}+\delta)\bigr\}.
\]
Using the well-known identity
\[
\alpha x^{3}+\beta x^{2}+\gamma x+\delta \;=\; \Sigma\,\frac{\alpha a^{3}+\beta a^{2}+\gamma a+\delta}{a-b\,.\,a-c\,.\,a-d}\,(b-x\,.\,c-x\,.\,d-x),
\]

% =====================================================================
% Page 208 of book (scan p. 228)
% =====================================================================
\newpage
\noindent\textsc{208 \hfill A Memoir on the Abelian and Theta Functions. \hfill [825}
\vspace{0.6em}

\noindent
where the summation extends to the four terms obtained by the cyclical interchanges
of the letters $a$, $b$, $c$, $d$: and the like identity for $\alpha x^{3}+\beta x^{2}+\gamma x+\delta$, there will be
terms in $\alpha a^{3}+\beta a^{2}+\gamma a+\delta$, $\alpha b^{3}+\beta b^{2}+\gamma b+\delta$, $\alpha c^{3}+\beta c^{2}+\gamma c+\delta$, but the term in
$\alpha d^{3}+\beta d^{2}+\gamma d+\delta$ will disappear of itself.  After some easy reductions, the result is
\[
\sqrt{\lambda}\,DE_{m}\,A_{m}B_{m}C_{m}
 \;=\; \Sigma\,\frac{\alpha a^{3}+\beta a^{2}+\gamma a+\delta}{b-a\,.\,c-a\,.\,d-a}\,B_{12}^{x}C_{12}^{x},
\]
where the summation extends to the three terms obtained by the cyclical interchanges
of the letters $a$, $b$, $c$: We have $\alpha a^{3}+\beta a^{2}+\gamma a+\delta = \sqrt{\mu}\,A_{a}A_{m}A_{m}$ and similarly for the
other two terms; the whole equation thus divides by $A_{m}B_{m}C_{m}$, and we find
\[
-\frac{\sqrt{\mu}}{\sqrt{\lambda}}\,(x_{1}-x_{3}\,.\,x_{1}-x_{4}\,.\,x_{2}-x_{3}\,.\,x_{2}-x_{4})\,DE_{m}
 \;=\; \frac{1}{x_{1}-x_{2}\,.\,x_{1}-x_{4}\,.\,x_{2}-x_{3}\,.\,x_{3}-x_{4}}\,M,
\]
where $M$ is a given rational and integral function of the single and double-letter
functions of $x_{1}$, $x_{2}$ and $x_{3}$, $x_{4}$.  The factor on the left-hand side has been made the
same as in the formula for the single-letter functions $A_{m}$, \&c., and to do this it was
necessary to bring in on the right-hand side the factor
\[
\frac{1}{x_{1}-x_{2}\,.\,x_{1}-x_{4}\,.\,x_{2}-x_{3}\,.\,x_{3}-x_{4}};
\]
this disappears in the expression for the ratio of two double-letter functions; but it
enters into the expression for the ratio of a single-letter to a double-letter function,
and it then requires to be itself expressed in terms of the functions of $x_{1}$, $x_{2}$ and
$x_{3}$, $x_{4}$: it is easy to see that we have
\[
\resizebox{\textwidth}{!}{$\displaystyle
\frac{1}{x_{1}-x_{2}\,.\,x_{1}-x_{4}\,.\,x_{2}-x_{3}\,.\,x_{3}-x_{4}}
 \;=\; \Sigma\,\frac{(A^{x}_{12}B_{12}^{x}-A^{x}_{34}B_{34}^{x})\,(A^{x}_{12}C_{12}^{x}-A^{x}_{34}C_{34}^{x})}{(a-b)(a-c)},
$}
\]
where the summation extends to the three terms obtained by the cyclical interchanges
of the letters $a$, $b$, $c$: these being a set of any three out of the six letters.

We have, in what precedes, obtained the expressions for the ratios of the 16
functions $A_{m},\ldots,F_{56},AB_{m},\ldots,DE_{m}$ in terms of the ratios of the like functions of $x_{1}$, $x_{2}$
and $x_{3}$, $x_{4}$.

% =====================================================================
% Page 209 of book (scan p. 229)  -- Chapter VII begins
% =====================================================================
\newpage
\noindent\textsc{825] \hfill A Memoir on the Abelian and Theta Functions. \hfill 209}
\vspace{0.8em}

\begin{center}
{\large\scshape Chapter VII.\quad The Functions $T$, $U$, $V$, $\Theta$.}
\end{center}
\vspace{0.4em}

\noindent
The present chapter is substantially a reproduction of C.\ and G.'s seventh section,
``Die Function $T_{\mu}^{\kappa}$,'' borrowing only from the next section the definition of the
theta-function; but for greater simplicity I consider for the most part, the case, fixed
curve a quartic, $n=4$, $p=3$.

\bigskip
\begin{center}
\textit{Integral Form of the Affected Theorem.\quad Art.\ Nos.\ 164 to 169.}
\end{center}

\textbf{164.}\quad Writing for shortness $\dfrac{(x,y,z)_{0}^{n-2}\,d\omega}{012}=d\Pi_{12}$, we are concerned with the integrals
$\displaystyle\int_{a'}^{a}d\Pi_{12}$ which present themselves in connexion with the affected theorem: the notation
is explained, Chap.~V.; $a$, $a'$ are points on the curve $f$; the variable may be any
parameter serving for the determination of the current point, and the integral, taken
from the value which belongs to the point $a'$ to the value which belongs to the
point $a$, is represented as above by means of the two points $a$, $a'$ as limits of the
integral.  It is assumed that the integral is a canonical integral having the limits and
the parametric points interchangeable, $\displaystyle\int_{a'}^{a}d\Pi_{12} = \int_{2}^{1}d\Pi_{aa'}$: see Chapter IV.

\medskip
\textbf{165.}\quad Writing for shortness
\[
\Bigl(\int_{a'}^{a}+\int_{b'}^{b}+\int_{c'}^{c}+\ldots\Bigr)d\Pi_{12}
   \;=\;\int\!\binom{a,\,b,\,c,\ldots}{a',\,b',\,c',\ldots}d\Pi_{12},
\]
then if $\phi$, $\psi$ are curves each of the order $m$, the former of them intersecting the
fixed curve $f$ in the points $a$, $b$, $c$, \ldots, and the latter of them intersecting the same
curve in the points $a'$, $b'$, $c'$,\ldots, and if $\phi_{1}$, $\phi_{2}$, $\psi_{1}$, $\psi_{2}$ are what the functions $\phi$, $\psi$
become on substituting therein in place of the current coordinates the values which
belong to the parametric points 1, 2 respectively; the theorem becomes
\[
\int\!\binom{a,\,b,\,c,\ldots}{a',\,b',\,c',\ldots}d\Pi_{12} \;=\; \log\frac{\phi_{1}\psi_{2}}{\phi_{2}\psi_{1}}.
\]
The superior limits may be interchanged in any manner, and so also the inferior
limits may be interchanged in any manner.  If a superior limit coincide with an
inferior limit, the two may be thus considered as belonging to an integral which
will then have the value $0$, and the coincident points may therefore be omitted from
the expression on the left-hand side: and so in the case of any number of coincidences.

\medskip
\textbf{166.}\quad If the intersections of the curves $\phi$, $\psi$ and the parametric points are
situate on a curve of the order $m$; then taking the equation of this curve to be
$\phi+\lambda\psi = 0$, we have simultaneously $\phi_{1}+\lambda\psi_{1}=0$, $\phi_{2}+\lambda\psi_{2}=0$; whence $\phi_{1}\psi_{2}=\phi_{2}\psi_{1}$, and
the logarithmic term disappears: viz.\ the theorem becomes
\[
\int\!\binom{a,\,b,\,c,\ldots}{a',\,b',\,c',\ldots}d\Pi_{12} \;=\; 0.
\]
\smallskip
\hfill C.~XII.\hfill 27

% =====================================================================
% Page 210 of book (scan p. 230)
% =====================================================================
\newpage
\noindent\textsc{210 \hfill A Memoir on the Abelian and Theta Functions. \hfill [825}
\vspace{0.6em}

\noindent
\textbf{167.}\quad Suppose that the curves $\phi$, $\psi$ are each of them a major curve, that is, a
curve of the order $n-2$ passing through the $\delta$ dps, and consequently besides meeting
the curve $f$ in $n(n-2)-2\delta$, $=2p+n-2$ points: the theorem is
\[
\int\!\binom{a,\,b,\,c,\ldots}{a',\,b',\,c',\ldots}d\Pi_{12} \;=\; \log\frac{\phi_{1}\psi_{2}}{\phi_{2}\psi_{1}},
\]
where the numbers of the superior and of the inferior points are each $=2p+n-2$.

\medskip
\textbf{168.}\quad Suppose further that the curves $\phi$, $\psi$, being major curves as above, pass
each of them through the $n-2$ residues of 1, 2; they besides meet in $(n-2)(n-3)$
points, viz.\ these are the $\delta$ dps and $(n-2)(n-3)-\delta$ variable points: these $(n-2)(n-3)$
points lie on a minor curve, that is, a curve of the order $n-3$ passing through the
dps; and the minor curve together with the parametric line $12$ make together a
major curve, passing through the intersections of $\phi$, $\psi$ and also through the parametric
points 1, 2: viz.\ these points and the intersections of $\phi$, $\psi$ are situate on a curve of
the order $n-2$; the logarithmic term thus vanishes.  The intersections of $\phi$ with the
fixed curve are the $\delta$ dps, the $n-2$ residues and $2p$ other points, say these are
$a$, $b$, $c$,\ldots, $a^{x}$, $b^{x}$, $c^{x}$,\ldots; similarly the curve $\psi$ meets the fixed curve in the $\delta$ dps, the
$n-2$ residues, and in $2p$ other points, say these are $d$, $e$, $f$,\ldots, $d^{x}$, $e^{x}$, $f^{x}$,\ldots: the
theorem is
\[
\int\!\binom{a,\,b,\,c,\ldots;\,a^{x},\,b^{x},\,c^{x},\ldots}{d,\,e,\,f,\ldots;\,d^{x},\,e^{x},\,f^{x},\ldots}d\Pi_{12} \;=\; 0,
\]
where there are $2p$ superior and inferior points respectively.

\medskip
\textbf{169.}\quad I introduce the definitions: a minor curve meets the fixed curve in the
dps and in $2p-2$ other points, called ``cominors'': a major curve passing through
the $n-2$ residues of the points 1, 2, meets the fixed curve in the $\delta$ dps, the $n-2$
residues and in $2p$ other points, called ``comajors in regard to the points 1, 2.''
Observe that $p-1$ of the cominors determine uniquely the remaining $p-1$ cominors;
and similarly $p$ of the comajors determine uniquely the remaining $p$ comajors.

The foregoing theorem thus is that the sum $\displaystyle\int\!(\;)\,d\Pi_{12} = 0$, when the superior
points and the inferior points are each of them a system of comajors in regard to
the parametric points 1, 2.

\bigskip
\begin{center}\textit{Fixed Curve a Quartic.\quad Art.\ No.\ 170.}\end{center}

\textbf{170.}\quad It would be easy to go on with the general form; but as already mentioned,
I prefer to consider the case, fixed curve a quartic, $n=4$, $p=3$.  A minor curve is
here a line meeting the quartic in 4 points, which are ``cominors''; the major curve
is a conic, and if this passes through the residues of 1, 2 it besides meets the
quartic in 6 points, which are ``comajors in regard to the points 1, 2.''  Two points
and their residues are cominors, but this is only by reason that $n-3=1$.

\bigskip
\begin{center}\textit{The Function $T$.\quad Art.\ No.\ 171.}\end{center}

\textbf{171.}\quad In conformity with C.\ and G., I introduce the functional symbol
\[
T_{12}\binom{a,\,b,\,c,\ldots}{a',\,b',\,c',\ldots} \;=\; \int\!\binom{a,\,b,\,c,\ldots}{a',\,b',\,c',\ldots}d\Pi_{12},
\]

% =====================================================================
% Page 211 of book (scan p. 231)
% =====================================================================
\newpage
\noindent\textsc{825] \hfill A Memoir on the Abelian and Theta Functions. \hfill 211}
\vspace{0.6em}

\noindent
so that $T$ denotes a function of the parametric points, and of the sets of superior
and inferior points respectively.  The foregoing theorem for the quartic thus is
\[
T_{12}\binom{a,\,b,\,c,\,a^{x},\,b^{x},\,c^{x}}{d,\,e,\,f,\,d^{x},\,e^{x},\,f^{x}} \;=\; 0.
\]
Observing that
\[
\int_{d}^{a}+\int_{d^{x}}^{a^{x}} \;=\; 2\!\int_{2}^{1}-\int_{a^{x}}^{d}+\int_{d^{x}}^{a},
\]
and so in other cases, this may be written
\[
T_{12}\binom{a,\,b,\,c}{d,\,e,\,f} \;=\; T_{12}\binom{a^{x},\,b^{x},\,c^{x}}{a,\,b,\,c}-T_{12}\binom{d,\,e,\,f}{d^{x},\,e^{x},\,f^{x}};
\]
and if, as a definition of $T_{12}(a,b,c)$, we write
\[
T_{12}(a,b,c) \;=\; T_{12}\binom{a,\,b,\,c}{a^{x},\,b^{x},\,c^{x}} - 2\!\int_{2}^{1},
\]
where $a^{x}$, $b^{x}$, $c^{x}$ are the comajors of $a$, $b$, $c$ in regard to 1, 2, then the equation is
\[
T_{12}\binom{a,\,b,\,c}{d,\,e,\,f} \;=\; T_{12}(a,b,c)-T_{12}(d,e,f),
\]
viz.\ the function of the $(2p+2=)8$ points $1$, $2$, $a$, $b$, $c$, $d$, $e$, $f$ is here expressed as
a difference of two functions each of $(p+2=)5$ points: $T_{12}(a,b,c)$ is regarded as a
function of the $5$ points $1$, $2$, $a$, $b$, $c$, because the remaining points $a^{x}$, $b^{x}$, $c^{x}$ depend
only on these $5$ points.

\bigskip
\begin{center}\textit{The Function $U$.\quad Art.\ Nos.\ 172 to 175.}\end{center}

\textbf{172.}\quad We consider on the quartic the points $\xi$, $\mu$; $1$, $2$, $3$; and take $f$, $f'$ for
the cominors of $2$, $3$; $g$, $g'$ for the cominors of $3$, $1$; and $h$, $h'$ for the cominors of
$1$, $2$.  We write
\[
T = T_{\mu}(1,2,3),
\]
it is to be shown that there exists a function $U(1,2,3;\xi)$ such that
\[
\delta U \;=\; \tfrac{1}{2}\bigl[\delta_{\xi}T + \delta_{1}T_{1} + \delta_{2}T_{2} + \delta_{3}T_{3}\bigr],
\]
viz.\ considering $\xi$, $1$, $2$, $3$ as variable points on the quartic, the whole infinitesimal
variation of $U$ is the sum of these parts, where $\delta_{\xi}T$ is the variation of $T$ when only
$\xi$ is varied, $\delta_{1}T_{1}$ the variation of $T_{1}$ when only $1$ is varied, and similarly for $\delta_{2}T_{2}$
and $\delta_{3}T_{3}$.  We consider in the proof three other points $4$, $5$, $6$ on the quartic; and
taking $l$, $l'$ for the cominors of $5$, $6$; $m$, $m'$ for those of $6$, $4$; and $n$, $n'$ for those
of $4$, $5$, we write further
\[
X_{1} = T_{\mu}(4,l,l'),\qquad X_{2} = T_{\mu}(5,m,m'),\qquad X_{3} = T_{\mu}(6,n,n').
\]
\smallskip
\hfill $27\,$---$\,2$

% =====================================================================
% Page 212 of book (scan p. 232)
% =====================================================================
\newpage
\noindent\textsc{212 \hfill A Memoir on the Abelian and Theta Functions. \hfill [825}
\vspace{0.6em}

\noindent
and it then requires to be shown that
\begin{align*}
\tfrac{1}{2}\,\delta T          &= \int\!\binom{123}{456}d\Pi_{1\mu}+\tfrac{1}{2}\,T_{1\mu}(4,5,6),\\
\tfrac{1}{2}(T_{1}-X_{1}) &= \int\!\binom{456}{423}d\Pi_{1\mu}+\log\frac{\mu23\,.\,156}{\mu56\,.\,123},\\
\tfrac{1}{2}(T_{2}-X_{2}) &= \int\!\binom{564}{531}d\Pi_{1\mu}+\log\frac{\mu31\,.\,264}{\mu64\,.\,231},\\
\tfrac{1}{2}(T_{3}-X_{3}) &= \int\!\binom{645}{612}d\Pi_{1\mu}+\log\frac{\mu12\,.\,345}{\mu45\,.\,312},
\end{align*}
where $\mu23$ is the determinant formed with the coordinates of the points $\mu$, $2$, $3$
respectively; and so in other cases.

\medskip
\textbf{173.}\quad We have
\[
T_{1\mu}\binom{1,\,2,\,3}{4,\,5,\,6} \;=\; \tfrac{1}{2}\,T_{1\mu}(1,2,3) - \tfrac{1}{2}\,T_{1\mu}(4,5,6),
\]
that is,
\[
=\;\tfrac{1}{2}\,T \;-\;\tfrac{1}{2}\,T_{1\mu}(4,5,6),
\]
and thence the above value of $\tfrac{1}{2}\,\delta T$.

The affected theorem gives
\[
\int\!\binom{f,\,f'\!,\,2,\,3}{l,\,l'\!,\,5,\,6}d\Pi_{1\mu} \;=\; \log\frac{F_{1}L_{\mu}}{F_{\mu}L_{1}},
\]
where $F=0$ is the equation of the line through $f$, $f'$, $2$, $3$; and $F_{1}$, $F_{\mu}$ are what the
function $F$ becomes, on substituting therein for the current coordinates the coordinates
of the points $1$, $\mu$ respectively.  And similarly $L=0$ is the equation of the line
through $l$, $l'$, $5$, $6$; and $L_{1}$, $L_{\mu}$ are what the function $L$ becomes by the same sub\-
stitutions respectively. The values of $F_{1}$, $F_{\mu}$ are $123$, $\mu23$: those of $L_{1}$, $L_{\mu}$ are $156$,
$\mu56$: and the logarithmic term is thus
\[
=\;\log\frac{\mu23\,.\,156}{\mu56\,.\,123}.
\]
We then have
\[
\tfrac{1}{2}(T_{1}-X_{1}) \;=\; \int\!\binom{f,\,f'\!,\,2,\,3}{4,\,l,\,l'} d\Pi_{1\mu}
   \;=\; \int\!\binom{f,\,f'\!,\,2,\,3}{l,\,l'\!,\,5,\,6}d\Pi_{1\mu},
\]
and in this last expression for $\tfrac{1}{2}(T_{1}-X_{1})$ substituting for the second term the
logarithmic value just obtained, we have the required expression for $\tfrac{1}{2}(T_{1}-X_{1})$: and
those for $\tfrac{1}{2}(T_{2}-X_{2})$ and $\tfrac{1}{2}(T_{3}-X_{3})$ are deduced by mere cyclical permutations of the
letters.

\medskip
\textbf{174.}\quad Returning to the assumed relation $\delta U = \tfrac{1}{2}\,[\delta_{\xi}T + \delta_{1}T_{1} + \delta_{2}T_{2} + \delta_{3}T_{3}]$, in order
to the existence of the function $U$, it is only necessary to show that $T-T_{1}$ contains
no term in $1$, $\xi$, that is, no term depending on both these points, and that
$T_{1}-T_{2}$ contains no term in $1$, $2$; for then, by symmetry, the like properties hold in
regard to $T-T_{2}$, $T-T_{3}$, $T_{1}-T_{3}$, $T_{2}-T_{3}$ respectively, and the assumed expression is a
complete differential, from which the function $U$ may be obtained by integration.

% =====================================================================
% Page 213 of book (scan p. 233)
% =====================================================================
\newpage
\noindent\textsc{825] \hfill A Memoir on the Abelian and Theta Functions. \hfill 213}
\vspace{0.6em}

\noindent
\textbf{175.}\quad To show that $T-T_{1}$ contains no term in $1$, $\xi$.

For $T$, the only term in $1$, $\xi$ is $\displaystyle\int_{1}^{\xi}d\Pi_{1\mu}$,

\hspace{4em}``\quad $T_{1}$\quad'' \quad ``\quad ``\quad\;\;\;$\displaystyle\int_{1}^{\xi}d\Pi_{1\mu}$,

\noindent and it is to be shown that the difference of the two integrals contains no term in
$1$, $\xi$.  Considering on the quartic the two new points $\delta$, $e$, the first integral is
\[
\int\!\binom{1,\,\xi}{\delta,\,e}\,(d\Pi_{1\mu}+d\Pi_{1\mu}) \;=\; \int_{\delta}^{1}d\Pi_{1\mu}+\int_{e}^{\xi}d\Pi_{1\mu}+\int_{1}^{\xi}d\Pi_{1\mu},
\]
and the second is
\[
\int\!\binom{\xi,\,e}{\delta,\,4}\,(d\Pi_{1\mu}+d\Pi_{1\mu}) \;=\; \int_{\delta}^{\xi}d\Pi_{1\mu}+\int_{4}^{e}d\Pi_{1\mu}+\int_{1}^{4}d\Pi_{1\mu}.
\]
Hence in the difference the only terms which can contain $1$, $\xi$ are
\[
\int_{e}^{1}d\Pi_{1\mu}+\int_{\xi}^{4}d\Pi_{1\mu},
\]
and this is $=0$: wherefore there is not in the difference any term in $1$, $\xi$.
This proves the property for $T-T_{1}$.  The property for $T_{1}-T_{2}$ is proved in a similar manner.

\bigskip
\begin{center}\textit{Theorems in regard to the Function $U$.\quad Art.\ Nos.\ 176 to 179.}\end{center}

\textbf{176.}\quad Theorem (A). To prove
\begin{equation*}
U(1,2,3;\xi) - U(1,2,3;\mu) \;=\; \tfrac{1}{2}\,T_{\xi\mu}(1,2,3),\tag{A}
\end{equation*}
we have
\begin{align*}
U(1,2,3;\xi)-U(1,2,3;\mu) &= \int_{\mu}^{\xi}d_{\xi}U = \tfrac{1}{2}\!\int_{\mu}^{\xi}d_{\xi}T\\
                          &= \tfrac{1}{2}\,T_{\mu}(1,2,3)-\tfrac{1}{2}\,T_{\mu}(1,2,3),
\end{align*}
and $T_{\xi\mu}(1,2,3) = \displaystyle\int\!\binom{1,\,2,\,3}{1^{x},\,2^{x},\,3^{x}}d\Pi_{\xi\mu}$, where $d\Pi_{\xi\mu} = 0$, viz.\ considering this as derived
from $d\Pi_{\xi\mu}$, $= Q_{\xi\mu}d\omega$, by making the point $\xi$ coincide with $\mu$, then when $\xi$ is indefinitely
near to $\mu$, the numerator and denominator of $Q_{\xi\mu}$ are each of them infinitesimal of
the orders $3$ and $2$ respectively, and thus the function $Q_{\xi\mu}$ ultimately vanishes (see as
to this, Chap.~V.\ Art.\ Nos.\ 99 to 106).  We have therefore $T_{\xi\mu}(1,2,3) = 0$, and the
required theorem is proved.

\medskip
\textbf{177.}\quad Theorem (B). To prove
\begin{equation*}
U(1,2,3;\xi) - U(4,2,3;\xi) \;=\; \tfrac{1}{2}\,T_{1\xi}(f,f'),\tag{B}
\end{equation*}
where, as before, $f$, $f'$ are the cominors of $2$, $3$, that is, $2$, $3$, $f$, $f'$ lie on a line, we have
\begin{align*}
U(1,2,3;\xi)-U(4,2,3;\xi) &= \int_{4}^{1}d_{1}U = \tfrac{1}{2}\!\int_{4}^{1}d_{1}T\\
                          &= \tfrac{1}{2}\,T_{\mu}(\xi,f,f')-\tfrac{1}{2}\,T_{4\mu}(\xi,f,f'),
\end{align*}

% =====================================================================
% Page 214 of book (scan p. 234)
% =====================================================================
\newpage
\noindent\textsc{214 \hfill A Memoir on the Abelian and Theta Functions. \hfill [825}
\vspace{0.6em}

\noindent
the point $\mu$ is arbitrary, and it may be taken to coincide with $4$; but we then have
$T_{4\mu}(\xi,f,f') = 0$, and the theorem is thus proved.

\medskip
\textbf{178.}\quad Theorem (C). We have
\begin{equation*}
T_{1\xi^{x}}(\xi,f,f') + T_{1\xi^{x}}(\eta,k,k') \;=\; 0;\tag{C}
\end{equation*}
where $\xi$, $\eta$, $1$, $2$, $3$ are arbitrary points on the quartic; $1^{x}$, $2^{x}$, $3^{x}$ are the comajors
of $1$, $2$, $3$ in regard to $\xi$, $\eta$, viz.\ the points $1$, $2$, $3$, $1^{x}$, $2^{x}$, $3^{x}$ lie on a conic which
passes through $\xi$, $\eta$: the residues (or cominors) of $\xi$, $\eta$: $f$, $f'$ are the cominors of
$2$, $3$; and $k$, $k'$ are the cominors of $2^{x}$, $3^{x}$.

Taking $\theta$, $\theta'$ for the cominors of $1$, $1^{x}$, the four lines $\xi\eta\theta\theta'$, $11^{x}\theta\theta'$, $23ff'$ and
$2^{x}3^{x}k k'$ form a quartic cutting the fixed quartic in the $16$ points: but of these,
$\xi$, $\eta$, $1$, $2$, $3$, $1^{x}$, $2^{x}$, $3^{x}$ lie in a conic: hence the remaining $8$ points $\theta$, $\theta'$, $\xi$, $\eta$,
$f$, $f'$, $k$, $k'$ lie in a conic: that is, $\xi$, $\eta$, $f$, $f'$, $k$, $k'$ lie on a conic through $\theta$, $\theta'$, the
residues of $1$, $1^{x}$: or they are comajors in regard to $1$, $1^{x}$; whence the theorem.

\medskip
\textbf{179.}\quad We have
\begin{center}
\textit{From A.}\hfill\textit{From B.}
\end{center}
\vspace{-1em}
\begin{align*}
\tfrac{1}{2}T_{1\xi^{x}}(\xi,f,f') &= U(f,f';1) - U(f,f';1^{x}) = U(1,2,3;\xi) - U(1^{x},2,3;\xi),\\
\tfrac{1}{2}T_{1\xi^{x}}(\eta,k,k') &= U(\eta,k,k';1) - U(\eta,k,k';1^{x}) = U(1,2^{x},3^{x};\eta) - U(1^{x},2,3;\eta);
\end{align*}
viz.\ we have thus two expressions for each term of the equation (C),
\[
T_{1\xi^{x}}(\xi,f,f') + T_{1\xi^{x}}(\eta,k,k') \;=\; 0.
\]

In particular, we have Theorem (D)
\begin{equation*}
U(1,2,3;\xi) - U(1^{x},2,3;\xi) \;=\; -U(1,2^{x},3^{x};\eta) + U(1^{x},2^{x},3^{x};\eta). \tag{D}
\end{equation*}

Again, we have $T_{1\xi}(1,2,3) + T_{1\xi}(1^{x},2^{x},3^{x}) = 0$; where $1$, $2$, $3$, $1^{x}$, $2^{x}$, $3^{x}$ are
comajors in regard to $\xi$, $\eta$:\ and
\begin{align*}
\tfrac{1}{2}\,T_{\xi\eta}(1,2,3)            &= U(1,2,3;\xi)-U(1,2,3;\eta),\\
\tfrac{1}{2}\,T_{\xi\eta}(1^{x},2^{x},3^{x}) &= U(1^{x},2^{x},3^{x};\xi)-U(1^{x},2^{x},3^{x};\eta);
\end{align*}
whence Theorem (E),
\begin{equation*}
U(1,2,3;\xi) - U(1,2,3;\eta) \;=\; -U(1^{x},2^{x},3^{x};\xi) + U(1^{x},2^{x},3^{x};\eta).\tag{E}
\end{equation*}

\bigskip
\begin{center}\textit{The Function $V$.\quad Art.\ Nos.\ 180 to 182.}\end{center}

\textbf{180.}\quad It is convenient to consider $U$ as a logarithm, say
\[
-U(1,2,3;\xi) \;=\; \log V(1,2,3;\xi),\qquad\text{or}\qquad V(1,2,3;\xi) = \exp.{-U(1,2,3;\xi)};
\]
$V$, like $U$, is a function of the $(p+1=)\,4$ points $1$, $2$, $3$, $\xi$, on the quartic.

The equation (D) thus becomes
\[
\frac{V(1,2,3;\xi)\;V(1^{x},2,3;\xi)}{V(1^{x},2^{x},3^{x};\eta)\;V(1,2^{x},3^{x};\eta)},
\]

% =====================================================================
% Page 215 of book (scan p. 235)
% =====================================================================
\newpage
\noindent\textsc{825] \hfill A Memoir on the Abelian and Theta Functions. \hfill 215}
\vspace{0.6em}

\noindent
where $1$, $2$, $3$, $1^{x}$, $2^{x}$, $3^{x}$ are comajors in regard to $\xi$, $\eta$: the equation shows that, in
the function $\dfrac{V(1,2,3;\xi)}{V(1^{x},2^{x},3^{x};\eta)}$, we can without alteration of the value interchange a
pair of points $1$, $1^{x}$ out of the system of comajor points; and it of course follows
that we can in any manner whatever interchange these points, so as to have any
three of them in the numerator-function and the remaining three in the denominator-function.
In particular, we have
\[
\frac{V(1,2,3;\xi)}{V(1^{x},2^{x},3^{x};\eta)} \;=\; \frac{V(1^{x},2^{x},3^{x};\xi)}{V(1,2,3;\eta)}.
\]
The equation (E) becomes
\[
\frac{V(1,2,3;\xi)\,V(1^{x},2^{x},3^{x};\eta)}{V(1,2,3;\eta)\,V(1^{x},2^{x},3^{x};\xi)},
\]
and multiplying we find
\[
V(1,2,3;\xi)^{2} \;=\; V(1^{x},2^{x},3^{x};\eta)^{2},
\]
that is, $V(1,2,3;\xi) = \pm V(1^{x},2^{x},3^{x};\eta)$, the sign being determinately $+$ or determ\-inately $-$,
according to the precise definition of the function $V$.

\medskip
\textbf{181.}\quad Considering $\xi$ as fixed points on the curve but $\eta$ as also; and $1^{x}$, $2^{x}$, $3^{x}$
a variable point (that is, the parametric line $\xi\eta$ as rotating about the fixed point $\eta$),
the points $1$, $2$, $3$ are then determined as the remaining intersections with the quartic of
the conic which passes through the points $1^{x}$, $2^{x}$, $3^{x}$ and through the points $\xi$, $\eta$, which
are the residues of $\xi$, $\eta$.  And by the theorem just obtained it appears that $1$, $2$, $3$,
thus varying, the function $V(1,2,3;\xi)$ remains constant.  This comes to saying that
$V$, considered as a function of the points $1$, $2$, $3$, $\xi$, satisfies a certain linear partial
differential equation of the first order, having a solution $V = F(u,v,w)$, an arbitrary
function of $u$, $v$, $w$, determinate functions of the points $1$, $2$, $3$, $\xi$.  And if we can
find $u$, $v$, $w$ functions of these points such that they each of them remain constant
when the points $1$, $2$, $3$, $\xi$ vary as above, then the arbitrary function of $u$, $v$, $w$ will
remain constant for the variation in question and will thus be a value of the function $V$.

\medskip
\textbf{182.}\quad It is easily seen that such functions are
\[
u,\;v,\;w \;=\; \Bigl(\int^{1}+\int^{2}+\int^{3}-\int^{\xi}\Bigr)\,x\,d\omega,\;\;y\,d\omega,\;\;z\,d\omega,
\]
the inferior limits being given points which are regarded as absolute constants.  For
by the pure theorem, we have
\[
\Sigma\,(x,y,z)_{1}\,d\omega \;=\; 0,
\]
where $(x,y,z)_{1}$ is an arbitrary linear function, and where the summation extends to
all the intersections of the quartic with any given curve.  Writing
\[
p \;=\; \int x\,d\omega,\;\;\int y\,d\omega\;\;\text{or}\;\;\int z\,d\omega,\quad\text{that is,}\quad p = \!\int^{1}\!x\,d\omega,\;\int^{1}\!y\,d\omega\;\text{or}\;\int^{1}\!z\,d\omega,
\]
the inferior limits being any absolutely fixed point on the curve, and similarly
$p_{2}$, \&c.; the integral form of the theorem is $\Sigma p = \text{constant}$.  And applying the theorem

% =====================================================================
% Page 216 of book (scan p. 236)
% =====================================================================
\newpage
\noindent\textsc{216 \hfill A Memoir on the Abelian and Theta Functions. \hfill [825}
\vspace{0.6em}

\noindent
successively to the parametric line, and to the conic which determines the points
$1$, $2$, $3$, we have
\begin{align*}
p_{\xi}+p_{\eta}+p_{\xi'}+p_{\eta'} &= \text{const.},\\
p_{\xi}+p_{\eta}+p_{1}+p_{2}+p_{3}+p_{1^{x}}+p_{2^{x}}+p_{3^{x}} &= \text{const.}
\end{align*}
Taking the difference of these equations, we have
\[
p_{1}+p_{2}+p_{3}-p_{\xi'}+p_{1^{x}}+p_{2^{x}}+p_{3^{x}}-p_{\eta'} \;=\; \text{const.},
\]
viz.\ the points $\eta$, $1^{x}$, $2^{x}$, $3^{x}$ being fixed points, this is
\[
p_{1}+p_{2}+p_{3}-p_{\xi} \;=\; \text{const.},
\]
that is, the functions $u$, $v$, $w$ defined as above are each of them constant under the
variation in question.

\bigskip
\begin{center}\textit{The Function $\Theta$.\quad Art.\ Nos.\ 183 and 184.}\end{center}

\textbf{183.}\quad The function $V(1,2,3;\xi)$ of the $(p+1=)\,4$ points $1$, $2$, $3$, $\xi$, is thus a
function of the $(p=)\,3$ arguments
\[
u,\;v,\;w \;=\; \Bigl(\int^{1}+\int^{2}+\int^{3}-\int^{\xi}\Bigr)\,x\,d\omega,\;\;y\,d\omega,\;\;z\,d\omega.
\]
Disregarding a constant and exponential factor, we say that it is a theta-function of
these arguments, and we write the result provisionally in the form
\[
V(1,2,3;\xi) \;=\; \Theta(u,v,w),
\]
the more precise definition of the theta-function being reserved for further con\-sideration.

\medskip
\textbf{184.}\quad It appears by what precedes that a sum of $(p=)\,3$ integrals
$\displaystyle\int\!\binom{abc}{def}d\Pi_{12}$,
otherwise called $T_{12}\binom{a,\,b,\,c}{d,\,e,\,f}$, is in the first place expressed (see No.\ 171) as a difference,
$=\tfrac{1}{2}\,T_{12}(a,b,c)-\tfrac{1}{2}\,T_{12}(d,e,f)$, of two functions $T$.  Each of these is by Theorem (A)
(No.\ 176) expressed as a difference of two functions $U$, that is, as the difference of
the logarithms, or logarithm of the quotient, of two functions $V$: such function $V$
according to its original definition a function of $(p+1=)\,4$ points, but in such wise
that the function is expressible as a function of $(p=)\,3$ arguments, and so expressed
it is a $\Theta$-function of these arguments: and the final result thus is that the sum of
$(p=)\,3$ integrals $\displaystyle\int\!\binom{a,b,c}{d,e,f}d\Pi_{12}$ is equal to the logarithm of a fraction, whereof the
numerator and the denominator are each of them a product of two $\Theta$-functions.

% =====================================================================
% Page 217 of book (scan p. 237) -- Paper 826
% =====================================================================
\newpage
\noindent\textsc{826] \hfill\hfill 217}
\vspace{1.4em}

\begin{center}
{\Large\bfseries 826.}\\[0.6em]
{\large\scshape Note on a Partition-Series.}\\[0.6em]
{\small [From the \textit{American Journal of Mathematics}, vol.\ vi.\ (1884), pp.\ 63, 64.]}
\end{center}
\vspace{0.6em}

\noindent
\textsc{Prof.\ Sylvester}, in his paper, ``A Constructive theory of Partitions, \&c.,'' \textit{American
Journal of Mathematics}, vol.\ v.\ (1883), p.\ 282, has given the following very beautiful
formula
\[
\resizebox{\textwidth}{!}{$\displaystyle
(1+ax)(1+ax^{2})(1+ax^{3})\ldots \;=\; 1 + \frac{(1+ax^{2})\,xa}{1-x} + \frac{1}{1-x\,.\,1-x^{2}}(1+ax)(1+ax^{3})\,x^{4}a^{2}
$}
\]
\[
\resizebox{\textwidth}{!}{$\displaystyle
\qquad\;\; +\;\frac{1}{1-x\,.\,1-x^{2}\,.\,1-x^{3}}(1+ax)(1+ax^{2})(1+ax^{4})\,x^{12}a^{3}+\ldots,
$}
\]
or, as this may be written,
\[
\Omega \;=\; 1 + P + Q(1+ax) + R(1+ax)(1+ax^{2}) + S(1+ax)(1+ax^{2})(1+ax^{3}) + \ldots
\]
where
\[
P = \frac{(1+ax^{2})\,xa}{\mathbf{1}},\qquad
Q = \frac{(1+ax^{4})\,x^{5}a^{2}}{\mathbf{1.2}},\qquad
R = \frac{(1+ax^{6})\,x^{12}a^{3}}{\mathbf{1.2.3}},\qquad
S = \frac{(1+ax^{8})\,x^{22}a^{4}}{\mathbf{1.2.3.4}},\;\;\&\text{c.,}
\]
the heavy figures $\mathbf{1}$, $\mathbf{2}$, $\mathbf{3}$, $\mathbf{4}$, \ldots\ of the denominators being, for shortness, written to
denote $1-x$, $1-x^{2}$, $1-x^{3}$, $1-x^{4}$, \ldots\ respectively.  The $x$-exponents $1$, $5$, $12$, $22$, \ldots\ are
the pentagonal numbers $\tfrac{1}{2}(3n^{2}-n)$.

To prove this, writing
\[
\resizebox{\textwidth}{!}{$\displaystyle
P' = \frac{ax^{2}}{\mathbf{1}},\quad
Q' = \frac{ax^{4}}{\mathbf{1}}+\frac{a^{2}x^{6}}{\mathbf{1.2}},\quad
R' = \frac{ax^{6}}{\mathbf{1}}+\frac{a^{2}x^{9}}{\mathbf{1.2}}+\frac{a^{3}x^{12}}{\mathbf{1.2.3}},\quad
S' = \frac{ax^{8}}{\mathbf{1}}+\frac{a^{2}x^{11}}{\mathbf{1.2}}+\frac{a^{3}x^{15}}{\mathbf{1.2.3}}+\frac{a^{4}x^{18}}{\mathbf{1.2.3.4}},\;\;\&\text{c.,}
$}
\]
where the $x$-exponents are
\[
2;\quad 3,\;3+4;\quad 4,\;4+5,\;4+5+6;\quad 5,\;5+6,\;5+6+7,\;5+6+7+8;\;\;\&\text{c.,}
\]
\smallskip
\hfill C.~XII.\hfill 28

\end{document}
